\section{Corpora and audit method}
\label{sec:method}

\subsection{Corpora}
\label{sec:corpora}
We audit seven corpora that release raw URL strings with binary labels (Table~\ref{tab:corpora}). We describe how each class was collected only where a publication we could check documents it, and otherwise report what the released strings show. Ebbu2017 pairs PhishTank URLs with URLs taken from Yandex search results \citep{sahingoz2019,rashid2024}. PhishStorm is a balanced set of 96,018 URLs released with the PhishStorm system \citep{marchal2014}. ISCX-URL2016 provides benign URLs crawled from top-ranked domains together with a phishing subset \citep{mamun2016,rashid2024}. Hannousse2021 is a balanced set of 11,430 URLs released with 87 page and URL features \citep{hannousse2021}. PhiUSIIL contains 235,795 URLs and is distributed by the UCI repository \citep{prasad2024}. Kaggle-MU is the benign and phishing part of the Kaggle malicious URL collection \citep{siddhartha2021}; its defacement and malware classes are dropped.

PT-Alexa is different in kind. It is a construction of ours, not a published corpus: it combines the PhishTank dump and the Alexa top-domain list that the GramBeddings repository ships in its PDRCNN folder \citep{bozkir2023,wang2019pdrcnn}, keeps every phishing URL and truncates the ranked legitimate list to the same size. The PDRCNN and GramBeddings papers used larger and unbalanced selections from the same two sources, so PT-Alexa should be read as a controlled example of the popularity-list design, not as a reproduction of either benchmark. We therefore lead with PhiUSIIL, a published corpus, whenever a result concerns the popularity-list design.

All files were obtained from public Git mirrors at fixed commits (Table~\ref{tab:prov}), so that every byte we analyse can be retrieved again; a script checks out each commit and verifies the SHA-256 digests of all \NSourceFiles{} source files. The original download pages of several corpora have moved or require registration, and mirrors are how many studies obtain these files. The PhiUSIIL file is the UCI archive itself. The PhishStorm and Kaggle-MU files are copies stored in the GramBeddings repository; for PhishStorm its reference file points to a third-party re-implementation of PhishStorm, which stores the training set, and for Kaggle-MU to the Kaggle page. Public copies of one corpus can differ in exactly the properties we audit; we found PhishStorm URLs with an explicit scheme in the mirror and without one inside Kaggle-MU. Every presentation statistic in this paper therefore describes the named mirror, and Section~\ref{sec:threats} returns to this limit.

\begin{table}[tp]
\centering
\caption{The seven corpora. Year: publication year of the corpus paper or release. Sources are listed only where a checked publication or the construction script documents them (--: not documented). Removed: exact duplicate strings plus strings carrying both labels. Unique: canonical strings with a single label, the unit used in all experiments; Raw $-$ Removed $-$ Unique is the number of canonical duplicates and canonical conflicts (Table~\ref{tab:overlap}). Phish: share of phishing URLs after exact de-duplication (\%).}
\label{tab:corpora}
\footnotesize
\begin{adjustbox}{max width=\linewidth}
\begin{tabular}{@{}lllp{2.4cm}rrrr@{}}
\toprule
Corpus & Year & Legitimate source & Phishing source & Raw & Removed & Unique & Phish \\
\midrule
Ebbu2017 \citep{sahingoz2019} & 2019 & Yandex search results & PhishTank & 73,575 & 18,812 & 51,721 & 45.4 \\
PhishStorm \citep{marchal2014} & 2014 & -- & -- & 96,018 & 0 & 85,598 & 50.0 \\
ISCX-URL2016 \citep{mamun2016} & 2016 & crawl of top domains & -- & 45,343 & 9 & 44,751 & 22.0 \\
PT-Alexa (this study) & constructed & Alexa top domains & PhishTank dump & 31,506 & 7 & 28,970 & 50.0 \\
Hannousse2021 \citep{hannousse2021} & 2021 & -- & -- & 11,430 & 1 & 11,070 & 50.0 \\
PhiUSIIL \citep{prasad2024} & 2024 & -- & -- & 235,795 & 425 & 231,706 & 42.7 \\
Kaggle-MU \citep{siddhartha2021} & 2021 & -- & -- & 522,214 & 54 & 509,092 & 18.0 \\
\bottomrule
\end{tabular}
\end{adjustbox}
\end{table}


\subsection{Three feature groups}
We separate the information in a URL into three groups. Each property was assigned to a group before any model was trained; the \texttt{https} flag was moved into a group of its own after the first internal review, for the reason given below.

\emph{Presentation} features describe how a string was recorded: an explicit scheme, a leading \texttt{www.} label, a slash at the end of the string, an HTML-escaped ampersand (\texttt{\&amp;}) and an upper-case letter in the host. A collector decides these properties when it writes a string down, and two strings that differ only in them usually lead a browser to the same page. They are not equivalent under the URI standard \citep{rfc3986}, which treats \texttt{www.} and a missing scheme as different references, so we make no equivalence claim; we claim only that in these corpora they mainly encode the collection pipeline. The split follows control: an attacker chooses the host (registered or compromised), path, query, fragment, userinfo and whether the server speaks \texttt{http} or \texttt{https}, while a collector chooses whether the scheme is written, whether \texttt{www.} is kept, whether a trailing slash or an escaped entity is stored, and whether a homepage or a deep page is sampled. Upper-case hosts occur in at most \UpperMaxPct\% of the URLs of any class and corpus.

\emph{Scheme type} is a single flag for an explicit \texttt{https} scheme. It is a real property of the server: phishing sites moved to HTTPS during the years these corpora span, and a collector that wrote out schemes records that change. In the PhishTank dump behind PT-Alexa, which records the URL a reporter submitted (often the lure link of a message, which is what a mail gateway sees, not necessarily the landing page), the share of submissions with an explicit \texttt{https} scheme rises from \PTHttpsEarly\% up to 2017 (\PTHttpsEarlyN{} submissions) to \PTHttpsTwenty\% in 2020 (\PTNTwenty). It is also a recording artefact when a curator adds a scheme that was not observed. We report every probe with and without this flag, and a canonical model with an added scheme-type token.

\emph{Structure} features describe the shape of the canonical URL: whether a path exists, path depth, whether a query exists, the number of query parameters and the canonical length. A corpus whose legitimate class consists of homepages differs from its phishing class in structure for sampling reasons.

\emph{Lexical} features describe the content of the canonical URL: an explicit non-default port, userinfo before an \texttt{@}, host length, dots, hyphens and digits in the host, an IP-address host, the subdomain count, the length of the public suffix, path length and digits, special characters, \texttt{@} anywhere, a double slash inside the path, percent-encoding, character entropy, a count of 14 brand and credential keywords, the longest token and the digit ratio. The 19 lexical features are used together with the five structure features.

\subsection{Canonicalisation and host view}
Canonicalisation removes the presentation features and keeps everything an attacker controls. It replaces the escaped ampersand \texttt{\&amp;} by \texttt{\&} and decodes no other entity, drops the scheme, lower-cases the host, strips trailing dots and one leading \texttt{www.} label from the host, and removes the default ports 80 and 443. Userinfo, other ports, the path, the query and the fragment are kept byte for byte, except that an empty query or fragment marker (a bare \texttt{?} or \texttt{\#}) is dropped. A single trailing slash is removed from the path only when no query or fragment follows. The canonicaliser is implemented without a general URL parser so that malformed strings keep their content, and 15 unit tests fix its behaviour on cases such as \texttt{http://user@www.paypal.com.evil.tk/x}, IPv6 hosts and query values that end in a slash.

The host view keeps only the canonical host. It removes presentation features and the path-sampling artefact of homepage-only legitimate sets, at the cost of genuine path information, and it serves as a diagnostic, not as a proposed detector. Registered domains (eTLD+1) are computed with the Public Suffix List snapshot bundled with \texttt{tldextract} \VerTldextract{}, with the private section enabled, so that \texttt{a.b.blogspot.com} and \texttt{c.blogspot.com} fall into different groups; IP hosts are their own group.

\subsection{Detectors}
Seven detectors are trained on every corpus. Four gradient-boosted tree models (scikit-learn \texttt{HistGradientBoostingClassifier}, 200 iterations, learning rate 0.1, 31 leaves) use the presentation flags, the presentation flags plus scheme type, the structure group and the lexical group. Three logistic-regression models use TF-IDF character n-grams ($n=2$ to $5$, minimum document frequency 2, at most 300,000 features, sublinear term frequency, case preserved, $C=1$) on the raw string, the canonical string and the host view. The hyper-parameters were fixed before the first run and were not tuned \citep{pedregosa2011}.

The character n-gram model is a simple stand-in for the neural detectors reviewed in Section~\ref{sec:related}. To check that the findings do not depend on that choice, we also train the character branch of URLNet \citep{le2018urlnet}: a 32-dimensional character embedding, convolutions of width 3, 4, 5 and 6 with 128 filters each, max-pooling, dropout 0.5 and a logistic output, trained for five epochs with Adam (learning rate $10^{-3}$, batch 256) on the first 200 characters. It is trained on raw and canonical strings under both within-corpus splits with three seeds.

\subsection{Protocols}
\paragraph{Within-corpus evaluation} Before sampling, every corpus is de-duplicated at the canonical level: rows that collapse to the same canonical string with the same label are merged, and canonical strings that carry both labels are dropped. The corpus is then subsampled to at most 25,000 URLs per class with a seed-specific random state and split 80/20 in two ways. The random split is stratified by label. The grouped split takes the first fold of a shuffled, stratified group five-fold partition with registered domains as groups, so that no registered domain occurs on both sides while the test share and the class prior stay close to those of the random split (split statistics are stored with the results). Ten seeds vary the sample and the split. For Hannousse2021 and PT-Alexa the cap exceeds the class sizes, so the seeds vary only the split.

\paragraph{Cross-corpus transfer} A detector trained on one corpus (at most 25,000 URLs per class) is tested on each other corpus (at most 5,000 URLs per class). Test URLs are purged against the full training corpus, not only the training sample, in two ways: the URL purge removes test URLs whose canonical form occurs anywhere in the de-duplicated training corpus, and the stricter domain purge also removes test URLs whose registered domain occurs there. A pair is skipped when purging leaves only one class; this happens for training on Kaggle-MU and testing on ISCX-URL2016, and for holding out ISCX-URL2016 in the leave-one-corpus-out protocol below, because every legitimate ISCX-URL2016 URL also occurs in Kaggle-MU. Three seeds vary both samples.

\paragraph{Leave-one-corpus-out} A detector is trained on the union of six corpora (at most 8,000 URLs per class from each) and tested on the held-out corpus, purged against all six training corpora with the same two rules. Three seeds are used.

\paragraph{Planned and added analyses} The presentation, structure, lexical and character detectors, both within-corpus splits and the cross-corpus protocol were planned before any result existed. The host view and leave-one-corpus-out were added after the first within-corpus results. The probe ablation, the path-subset probe, the scheme-type flag, the domain purge, the stratified grouped split and the character CNN were added in response to a first internal review; the homepage subset, the canonical model with a scheme-type token, the Kaggle-MU-clean subset, and the concentration analysis of transfer gains were added after a second; thresholds chosen on a validation split within and across corpora, a second Kaggle-MU subset and a time split of the PhishTank dump were added after a third; the feature-attribution analysis (Section~\ref{sec:shap}) was added last and checked by the same three referees, who asked for the symmetric n-gram share and the same-model intervention. The experiments of the first two reviews were rerun from empty result folders after the second review; the third-review controls ran afterwards as a separate script on the same files (Section~\ref{sec:repro}).

\subsection{Metrics and statistics}
Within corpora the primary metric is the Matthews correlation coefficient (MCC) at a 0.5 probability threshold, which uses all four cells of the confusion matrix and is informative under class imbalance \citep{chicco2020}. Across corpora the primary metric is ROC-AUC, because a threshold fitted on one corpus is not calibrated for another and MCC at 0.5 then mixes ranking quality with calibration; MCC is reported next to it. Because every corpus is near balanced while phishing is rare in deployment, we also convert true- and false-positive rates into the precision a detector would have at 10\%, 1\% and 0.1\% phishing prevalence. These are round values, not estimates: 10\% stands for a pre-filtered stream such as user reports and 0.1\% for unfiltered traffic, and the rate of a real deployment is not known for these corpora. A threshold of 0.5 is rarely what a deployment uses, so for the character models we also hold out 20\% of the training part as validation data, fit the model on the rest, choose the threshold with the highest validation true-positive rate at a validation false-positive rate of at most 0.1\% (and 1\%), and report the realised rates on the test part and on other corpora.

Means are reported with 95\% Student-$t$ intervals over seeds. Repeated random splits share training data, so the ordinary paired $t$-test and the Wilcoxon test overstate the evidence; we compare two settings on the same splits with the corrected resampled $t$-test of \citet{nadeau2003}, which inflates the variance by the test-to-training ratio, and adjust the seven corpus-level $p$-values of each comparison and metric with Holm's procedure \citep{holm1979,demsar2006}; we do not adjust across the different comparisons, each of which answers its own question. For transfer, the 41 evaluable ordered pairs share training and test corpora, so they are not independent. We average each difference over the test corpora of a training corpus and test the seven corpus-level means with the Wilcoxon signed-rank test \citep{wilcoxon1945}, whose smallest attainable two-sided $p$ with seven units is 0.016 (rounded). Leave-one-corpus-out results are reported descriptively.

\subsection{Reproducibility}
\label{sec:repro}
The pipeline runs from one script in a fixed order: unit tests of the canonicaliser, data loading with provenance hashes, the audit, within-corpus experiments, probe ablations, the control subsets, cross-corpus experiments, n-gram inspection, aggregation, and table and figure generation; the CNN runs as a second script, the third-review controls as a third and the attribution analysis as a fourth; a final script repeats the descriptive audit and the aggregation and then stores a SHA-256 digest of every code file. Every experimental result in the text, tables and figures is generated from the per-run result files; none is typed by hand. The run took \RunMinutes{} minutes for the linear and tree models, \CNNHours{} hours for the CNN, \RThreeMinutes{} minutes for the third-review controls and \ShapMinutes{} minutes for the attribution analysis on four cores of an \VerCpu{} virtual machine; the logs of all scripts are stored with the results. SHAP \VerShap{} computes the attributions. Software versions are Python \VerPython{}, scikit-learn \VerSklearn{}, NumPy \VerNumpy{}, pandas \VerPandas{}, SciPy \VerScipy{}, tldextract \VerTldextract{} and PyTorch \VerTorch{}. The code, per-run results and manuscript source will be released in a public repository on acceptance, with a script that fetches the mirrors at the listed commits and checks their digests; Table~\ref{tab:prov} lists the data sources.

\begin{table}[tp]
\centering
\caption{Source mirrors, commits and the corpora taken from each. The PhiUSIIL file is a copy of the UCI archive stored in the listed repository; the PhishStorm and Kaggle-MU files are copies inside GramBeddings, whose reference files point to a PhishStorm re-implementation and to the Kaggle page; SHA-256 digests of all source files are in \texttt{results/provenance.json}.}
\label{tab:prov}
\footnotesize
\begin{tabular}{@{}>{\raggedright\arraybackslash}p{6.4cm}lp{3.6cm}@{}}
\toprule
Repository (github.com/) & Commit & Corpora \\
\midrule
\url{ebubekirbbr/pdd} & \texttt{90b394d8f9fb} & Ebbu2017 \\
\url{fcdalgic/GramBeddings} & \texttt{222157e1a61a} & PhishStorm, ISCX-URL2016, PT-Alexa, Kaggle-MU \\
\url{AbdulHamidNoubaa/Phishing-Detection-Using-PhiUSIIL-Phishing-URL-Dataset-Website-} & \texttt{2d9a1f00a36b} & PhiUSIIL \\
\url{TetianaYaskova/Web-page-Phishing-Detection-Dataset} & \texttt{9ba531ffc919} & Hannousse2021 \\
\bottomrule
\end{tabular}
\end{table}

